% #############################################################################
% This is Chapter 6
% !TEX root = ../main.tex
% #############################################################################
% Change the Name of the Chapter i the following line
\fancychapter{Conclusion}
% The following line allows to ref this chapter
\label{chap:conclusion}

This chapter brings together the main findings from the identification and
rewriting experiments. Section~\ref{sec:conclusion-summary} summarizes the main
empirical conclusions and their implications for Portuguese variety adaptation,
while Section~\ref{sec:conclusion-future-work} outlines future work on data,
evaluation, and staged training.

\section{Summary}
\label{sec:conclusion-summary}

{\tolerance=1500\looseness=-1
This work studied European and Brazilian Portuguese through variety
identification and bidirectional rewriting in a shared T5Gemma 2 encoder-decoder
framework. Lower-cost experiments with the 270M-parameter model were used to
select the supervised protocol and examine reward-design choices, while the
4B-parameter model was used for the main comparisons of control formulation,
supervision composition, and staged adaptation. The protocol-selection
experiments improved individual metrics under some variants, but none
consistently improved performance across both tasks and evaluation resources.\par}

{\tolerance=1500\looseness=-1
The clearest result from the 4B-parameter experiments concerns supervision
composition. GPT-Wikipedia produces substantially better Golden Collection
rewriting than the GPT-Wikipedia + FRMT data mix, while the effect of adding FRMT
on FRMT itself depends on rewriting direction and control formulation. The
identification task is less sensitive to this change. Adding FRMT improves FRMT
identification under both formulations, while its effect on the Golden
Collection is mixed. The performance differences between encoder and decoder
control also depend on rewriting direction. Decoder control performs better in
every BP-to-EP rewriting comparison, whereas EP-to-BP includes cases where
encoder control performs better. Removing identification supervision from the
encoder-controlled models does not improve rewriting overall.\par}

{\tolerance=1500\looseness=-1
The reported models also perform better numerically than the closest published
systems used for comparison. The encoder-controlled GPT-Wikipedia model exceeds
the published FRMT identification result reported by
\cite{sousa2025enhancingportuguesevarietyidentification}, while the
decoder-controlled GPT-Wikipedia model exceeds the strongest trained BP-to-EP
Golden Collection result reported by
\cite{sanches2024brazilianportugueseeuropeanportuguese}. These comparisons remain
contextual, since the previous systems use different training setups and
objectives.\par}

{\tolerance=1500\looseness=-1
The training stages have different effects depending on the supervision
condition. Stage A produces its largest gains when GPT-Wikipedia is used for
later specialization, particularly on the Golden Collection. Stage C changes
the GPT-Wikipedia conditions only modestly. However, after GPT-Wikipedia + FRMT
specialization, Stage C recovers a substantial part of the Golden Collection
performance lost at Stage B, while its effect on FRMT rewriting remains mixed.
The identification results remain close to the corresponding Stage B values
after Stage C despite the continuation stage using rewriting-only data.\par}

{\tolerance=1500\looseness=-1
Many Golden Collection references require little or no change, making
conservative outputs highly competitive. Models trained with GPT-Wikipedia
preserve the source much more often than models trained with the GPT-Wikipedia +
FRMT data mix. Adding FRMT can reduce under-editing when a change is required,
but it also increases broader reformulation. Inspection of the aligned data
shows that some FRMT pairs contain lexical substitutions, omissions, and other
changes in expressed information that go beyond the localized changes required
for EP/BP rewriting. Similar reference-side differences also occur in some
Golden Collection examples.\par}

{\tolerance=1500\looseness=-1
These reference properties also affect what can be concluded from automatic
evaluation. BLEU and TER measure agreement with a particular reference, but
that reference does not always isolate the variety-specific change required by
the task. The same issue affects Stage C when sentence-level BLEU and TER
contribute to the reward, since a broader-than-required reference can favor
candidates that reproduce that broader rewrite. Across the experiments,
supervision composition produces the largest changes in rewriting behavior,
while the error analysis shows that the construction of the aligned pairs
influences both the behavior learned by the models and how that behavior is
evaluated.\par}

\section{Future Work}
\label{sec:conclusion-future-work}

{\tolerance=1500\looseness=-1
The strongest direction for future work is to improve the supervision used for
controlled rewriting. Larger manually verified resources would make it possible
to distinguish required variety-specific changes from optional paraphrase and
changes in expressed information more reliably. The OpenSubtitles filtering
procedure could also be evaluated more systematically by varying the SimAlign
coverage threshold and allowed sentence-length ratio, and by comparing the
current XLM-R, word-level intersection setup with alternative encoders,
alignment granularities, and extraction methods. For Stage C data selection, the
existing structural-similarity and edit criteria could be supplemented with
content-preservation checks for entities, numbers, parenthetical information,
units, and other source--reference differences. This could remove pairs with
differences beyond the requested variety change despite high structural
similarity.\par}

{\tolerance=1500\looseness=-1
Improved supervision should be accompanied by evaluation criteria that are less
dependent on a single reference. Future evaluation should assess whether the
required variety-specific changes were made while preserving the source meaning
and avoiding unnecessary rewriting. Human evaluation would provide the
strongest basis for validating these judgments, while targeted linguistic
checks, multiple acceptable references, or learned evaluators could support
larger-scale evaluation.\par}

{\tolerance=1500\looseness=-1
A validated evaluation signal would also make it possible to study Stage C
under better controlled conditions. The current Stage C comparison changes both
the optimization procedure and the training subset, so a follow-up should vary
these factors independently. One comparison could keep the Stage C training
subset fixed and contrast supervised continuation with reward-based
optimization, while another could hold the optimization method fixed and vary
the training subset. Once these effects are separated, reward design could be
studied more directly by comparing individual quality and conservatism
components, explicitly including the unchanged source among the candidates,
varying candidate-group construction, or assigning feedback to the particular
spans that require adaptation. These experiments could also test whether Stage C identification remains stable
under stronger rewriting-focused optimization.\par}

{\tolerance=1500\looseness=-1
Finally, the stability of these findings should be tested beyond the present
resources and T5Gemma 2 configurations. Additional domains, model families, and
larger manually evaluated test sets would show whether the observed effects of
supervision composition, control formulation, and staged adaptation generalize
beyond the current experimental setting. This broader validation would be most
informative after supervision and evaluation are better controlled, since it
would help separate model-dependent differences from effects caused by the
available training and evaluation resources. The descriptive external
comparison in \Cref{chap:appendix-external-model-comparison} could also be
strengthened by evaluating a current frontier proprietary model under the same
identification and bidirectional rewriting protocol.\par}
